
Three research questions are addressed:

\textbf{RQ1:} Does the structural efficiency measurement framework execute correctly end-to-end on real model checkpoints?

\textbf{RQ2:} Do GRPO and DPO exhibit different Semantic Edit Proportions under KL-matched conditions?

\textbf{RQ3:} What confounds affect checkpoint-comparison studies of RL fine-tuning, and how can they be detected?

\subsubsection{4.1 Datasets}

\begin{table}[htbp]
\centering
\begin{tabular}{llll}
\toprule
Dataset & Problems & Language & Source \\
\midrule
HumanEval+ & 164 & Python & evalplus [Liu et al., 2023] \\
MBPP+ & 378 & Python & evalplus [Liu et al., 2023] \\
\bottomrule
\end{tabular}
\end{table}

HumanEval+ and MBPP+ are used because they cover standard function-level Python code generation with explicit input/output contracts, making AST analysis well-defined. The evalplus harness provides reliable execution-based evaluation.

\subsubsection{4.2 Model}

\textbf{Base model:} DeepSeek-Coder-7B-instruct-v1.5 (\texttt{deepseek-ai/deepseek-coder-7b-instruct-v1.5}, HuggingFace), a 7B parameter decoder-only transformer specialized for code. The instruct-tuned variant is used as the starting point for post-training, consistent with standard practice \citep{Lambertetal2024}. Experiments run on NVIDIA H100 NVL (95,830 MiB), single GPU per run.

\subsubsection{4.3 Alignment Conditions}

\begin{table}[htbp]
\centering
\begin{tabular}{lll}
\toprule
Condition & Method & Reward Signal \\
\midrule
GRPO-binary & TRL GRPOTrainer & +1.0 (pass) / 0.0 (fail) \\
GRPO-error-type & TRL GRPOTrainer & Error taxonomy reward (SyntaxError=0.1, RuntimeError=0.3, AssertionError=0.7, pass=1.0) \\
DPO & TRL DPOTrainer & Execution-oracle preference pairs (intended; see Section 6.2, L4) \\
\bottomrule
\end{tabular}
\end{table}

Control variables held constant across conditions: base model, training problems (CodeAlpaca/OSS-Instruct subset), compute budget (1000 training steps), and evaluation harness.

\subsubsection{4.4 Sub-Experiments}

\textbf{h-e1 (Proof-of-Concept, RQ1):} A smoke test to verify measurement infrastructure. Executed on 6 HumanEval+ problems with hand-crafted code completions constructed to have known structural properties — GRPO completions with more control-flow and data-flow nodes, DPO completions with fewer. The purpose was infrastructure verification, not empirical measurement of real GRPO/DPO training behavior.

\textbf{h-m1 (Mechanism Analysis, RQ2):} Full structural analysis on KL-matched checkpoint pairs using real checkpoints. Designed for 27 matched pairs from 10+ distinct checkpoints. Due to checkpoint aliasing (Section 5.3), effective sample size was n_eff ≈ 2.

\subsubsection{4.5 Evaluation Metrics}

\textbf{Primary (framework diagnostic):}
\begin{itemize}
\item Semantic Edit Proportion (SEP): fraction of edits targeting CF+DF nodes
\item Semantic edit distance: ZSS edit distance restricted to CF+DF nodes
\item Structural efficiency: semantic edit distance per unit KL divergence
\end{itemize}

\textbf{Statistical tests:}
\begin{itemize}
\item Mann-Whitney U test on SEP distributions (α = 0.05)
\item Bootstrap 95\% CI on mean SEP differential (10,000 samples)
\end{itemize}

\textbf{Secondary (training confirmation):}
\begin{itemize}
\item Pass@1 on HumanEval+ and MBPP+ — used only to confirm training progressed, not as the primary comparison metric.
\end{itemize}
